\chapter{Visualisations and Frontend}\label{ch:viz}

Two visualisation approaches were developed to present the model's predictions in an accessible form.

\section{xT Heatmap}

For each version of the model, a heatmap of the predicted xT value across the pitch is generated by querying the model at every cell of a spatial grid, sweeping the ball position across all positions. The resulting probability values are plotted as a filled contour map using the \texttt{magma} colourmap (dark/purple = low threat, orange/bright = high threat), with standard pitch markings overlaid in white for reference.

For Version 1, a single heatmap is produced using a Pass action type at a fixed neutral game state. For Versions 2 and 3, the model is run separately for each of the four action types (Shot, Pass, Carry, Dribble), and a heatmap is produced for each action independently. Each grid cell therefore shows the model's predicted xT for a ball carrier at that position given that the specified action is taken. End positions are inferred per action type as follows: a shot's end position is fixed at the goal centre $(120, 40)$; a carry's end position is inferred by the start position projected toward goal by the empirical median carry distance of 4.06\,units ($\approx$\,4\,m on the standard 120$\times$80 pitch; see Table~\ref{tab:carry_stats}), computed from 217,606 carry events across the 292 training matches using the \texttt{carry.end\_location} field provided by StatsBomb; for a dribble the start position is used as both start and end, reflecting that dribbles are on-ball challenges without a directed ball displacement; for a pass xT is evaluated for each visible outfield teammate as a potential target, and the maximum value is reported.

Game-state assumptions.\ Because the heatmaps are generated without a live match context, the match-state scalar features are fixed at a neutral baseline throughout: period 1 (\texttt{period\_norm} $= 0.25$), minute 0 (\texttt{minute\_norm} $= 0.0$), and a tied scoreline (\texttt{score\_diff\_norm} $= 0.5$). The maps therefore represent positional threat under neutral game-state conditions, not a specific match situation.

A scenario comparison grid is also generated for Versions 2 and 3, showing how different defensive formations affect the threat landscape under the same neutral game-state assumptions. The scenarios tested are: no players (baseline), compact low block, high press, 2v1 overload in box, counter-attack 3v2, and goalkeeper off the line. These are discussed in detail in Chapters~\ref{ch:v2} and \ref{ch:v3}.

\section{Interactive Web Frontend}

To make the Version 3 model accessible and easy to explore, a Flask-based interactive web application was developed. The visual design and layout of the frontend interface were generated with the assistance of an AI coding tool. The frontend renders a football pitch as an HTML5 canvas and allows the user to:

\begin{itemize}
    \item Click to place the ball start position
    \item Place outfield teammates, outfield opponents, a friendly goalkeeper, and an opposing goalkeeper using mode buttons
    \item Click ``Predict'' to query the model in real time
\end{itemize}

The user first selects an action type (Shot, Pass, Carry, or Dribble) using the action selector. The Flask backend then receives the ball position, selected action type, and player positions as a JSON payload. It runs the model once for each of the four action types, inferring end positions from the player layout as described above, and returns the xT score for the selected action as the headline figure, alongside the scores for all four actions for comparison.

 As with the static heatmaps, match-state features are fixed at a neutral baseline (period 1, minute 0, score tied). The under pressure flag is inferred directly from player positions: if any goalside opponent is within 5\,m of the ball carrier, the flag is set to true.

The result is displayed as a single xT value for the chosen action type, with a per-action table showing all four scores simultaneously. This makes the model's output unambiguous: the displayed value is the raw model probability for one specific action from the current position, conditioned on the placed player layout and the neutral game-state baseline.

The screenshot (Figure~\ref{fig:frontend}) illustrates a representative inference scenario with Shot selected. The ball carrier is placed near the edge of the penalty area with teammates and opponents in the box. No pressure is detected (no opponent within 5\,m of the ball carrier).

The per-action table shows all four scores: Shot xT is 0.73044, reflecting genuine scoring danger from this position with the lane partially open; Pass xT reflects the quality of the best available teammate target; Carry xT reflects the value of advancing 4.06\,units ($\approx$\,4\,m) toward goal; and Dribble xT is low because no immediate pressing defender is present. Switching the selector to a different action type instantly updates the headline figure to the corresponding model score, allowing direct comparison of tactical options from any position.
